% Source: https://artofproblemsolving.com/community/c3h2355341_25_fun_problems
% Stable identity: aops-topic-2355341-q25
Let the bottom base be $c>8$. Dropping perpendiculars from the ends of the top base cuts off two right triangles with hypotenuse $5$ and horizontal leg $\dfrac{c-8}{2}$. A trapezoid exists exactly when $0<\dfrac{c-8}{2}<5$, i.e. $8<c<18$, so the greatest integer value is $c=17$.

\begin{asy}
usepackage("amsmath");
size(7cm);
defaultpen(fontsize(10pt));
real hh=sqrt(19)/2;
pair A=(0,0), B=(17,0), C=(12.5,hh), D=(4.5,hh), H1=(4.5,0), H2=(12.5,0);
draw(A--B--C--D--cycle,linewidth(0.9));
draw(D--H1,dashed); draw(C--H2,dashed);
draw((4.5,0.35)--(4.85,0.35)--(4.85,0)); draw((12.5,0.35)--(12.15,0.35)--(12.15,0));
label("$8$",(D+C)/2,N); label("$5$",(A+D)/2,NW); label("$5$",(B+C)/2,NE);
label("$4.5$",(A+H1)/2,S); label("$8$",(H1+H2)/2,S); label("$4.5$",(H2+B)/2,S);
label("$h$",(4.5,hh/2),E);
\end{asy}

Then the horizontal legs are $4.5$ and the height is
\[h=\sqrt{5^2-4.5^2}=\sqrt{25-20.25}=\sqrt{4.75}=\dfrac{\sqrt{19}}{2}.\]
The area is
\[\dfrac{8+17}{2}\cdot\dfrac{\sqrt{19}}{2}=\dfrac{25\sqrt{19}}{4},\]
choice $\boxed{\textbf{(D)}\ \dfrac{25\sqrt{19}}{4}}$.
